This companion partially formalizes Haar measure on actual topological groups with their Borel measurable structure. The shared module is \texttt{TraceFormulas.MeasureTheory.HaarBasics}. The manuscript assumes local compactness, the Hausdorff property and a countable base throughout; several library transport lemmas are valid under weaker hypotheses, so Lean omits unused assumptions for those results. This strengthens those statements without weakening the book's goal.

The book defines left Haar measure as nonzero, Radon and left invariant. Mathlib's \texttt{IsHaarMeasure} packages left invariance, finiteness on compact sets and positive measure of nonempty open sets. The shared bridge derives this positivity from nonzeroness and the stated regularity/topology. In the locally compact second-countable setting, the library's regularity instances also give compact inner approximation on every Borel set, as required by the book's Radon convention.

The declaration \texttt{haar\_positive\_on\_open} covers the positive-open-set lemma. The declaration \texttt{haar\_inverse} covers nonzeroness, regularity and right invariance of the inverse measure. It does not yet export the right-to-left and involution assertions as separate companion results. The declaration \texttt{haar\_unique\_scale} proves uniqueness of a finite strictly positive scalar for two left Haar measures. The proof uses Mathlib's invariant-measure uniqueness theorem, normalized on a compact set with nonempty interior. The prose instead proves uniqueness by a double-integral identity and compact-support approximation. These are different proof routes to the same left-Haar statement; right-Haar uniqueness still requires the corresponding exported transport argument.

The declaration \texttt{haar\_right\_translate\_scale} proves existence and uniqueness of the positive scalar for the measure obtained by pushing $\nu$ forward under $x\mapsto xh^{-1}$. The separate \texttt{haar\_right\_translate\_apply} proves that this pushforward evaluates on a Borel set $E$ as $\nu(Eh)$. This explicitly checks the convention in the \hyperref[eq:haar-modular-set-law]{modular set-transformation formula} \eqref{eq:haar-modular-set-law}.

The scalar's independence of normalization, its homomorphism and continuity, the integral transformation law, compact probability normalization, the abelian corollary, and the affine-group example remain unformalized in this companion. No construction of Haar measure on an arbitrary locally compact group is claimed: the prose assumes a Haar measure for these general results and constructs the measures needed in Chapter 1 explicitly. Machine checks and Steven's correspondence review are separate, and the latter remains not-reviewed.
